\section{Primitive-recursive fuel for fixed-dimensional lexicographic descent}
\label{sec:lex-fuel}
\label{sec:descent}

This section addresses a recurring gap between an ordinary termination
argument and an argument usable in intuitionistic primitive recursive
arithmetic.  A process may have a decreasing pair or a decreasing finite
tuple of natural numbers, while a lower coordinate is reset to a much larger
value whenever a higher coordinate decreases.  Its termination is familiar,
but neither an induction on an existential termination assertion nor an
unbounded while-loop is automatically available in iPRA.  We give an
explicit primitive-recursive fuel function.  The mathematical ingredients
are standard fixed-depth nested recursion (compare the broader length-function analysis in~\cite{schmitz2014}); the useful point is the fully
specified interface and the automatic construction of its growth envelope.

Throughout this section, the dimension $d$ is a fixed external natural
number.  It is part of a theorem instance, not an unrestricted input to the
resulting primitive-recursive program.  Tuple operations, finite sequence
operations, bounded maxima, bounded searches, and finite iteration are
primitive recursive.  Their bounded assertions are represented by their
primitive-recursive Boolean characteristic functions, so the inductions
below have quantifier-free representatives in iPRA.

\subsection{A bound on controlled descending chains}

Write $v>_{\mathrm{lex}}w$ when, at the first coordinate at which the two
$d$-tuples differ, the coordinate of $v$ is larger.  The first coordinate
is the most significant one.  There are no strict comparisons between
empty tuples.

Let $h(x,t)$ be a primitive-recursive function, where $x$ denotes an
arbitrary finite tuple of numerical parameters.  Suppose that iPRA proves
that $h(x,t)$ is nondecreasing in $t$.  Define, separately for each fixed
$d$, the following primitive-recursive functions:
\begin{align}
 F_0(x,s)&=s+1,\\
 F_{d+1}(x,s)&=F_d(x,\cdot)^{[h(x,s)+1]}(s).
 \label{eq:lex-fuel-recursion}
\end{align}
Here $f^{[r]}(s)$ means $r$ successive applications of $f$ to $s$, with
$f^{[0]}(s)=s$.  Formula~\eqref{eq:lex-fuel-recursion} describes a sequence
of primitive-recursive definitions.  It does not assert that the
three-variable function $(d,x,s)\mapsto F_d(x,s)$ is primitive recursive.

\begin{lemma}[Monotonicity of the bounds]
For every fixed $d$, iPRA proves
\[
 F_d(x,s)\ge s+1,
 \qquad
 s\le t\ \Longrightarrow\ F_d(x,s)\le F_d(x,t).
\]
\end{lemma}

\begin{proof}
The assertion for $d=0$ is immediate.  Suppose it has been proved for
$d$.  Iterating a nondecreasing function preserves inequalities between
starting points.  Iterating a function $f(u)\ge u+1$ more times cannot
decrease its value, and one iteration already gives a value at least one
larger than its starting point.  These two assertions follow by
quantifier-free induction on the iteration counter.  If $s\le t$, then
$h(x,s)+1\le h(x,t)+1$, and hence
\[
 F_d(x,\cdot)^{[h(x,s)+1]}(s)
 \le F_d(x,\cdot)^{[h(x,s)+1]}(t)
 \le F_d(x,\cdot)^{[h(x,t)+1]}(t).
\]
There is at least one iteration in the definition of $F_{d+1}$, so its
successor bound follows as well.  The induction on $d$ here constructs a
separate finite iPRA proof for each fixed dimension.
\end{proof}

\begin{theorem}[Controlled lexicographic chain bound]
\label{thm:controlled-lex}
Fix $d$.  Let
\[
 v_s,v_{s+1},\ldots,v_{s+L-1}\in\mathbb N^d
\]
be a finite sequence such that $v_t>_{\mathrm{lex}}v_{t+1}$ whenever both
terms occur, and every coordinate of $v_t$ is at most $h(x,t)$.  Then
\[
                         s+L\le F_d(x,s).
\]
For each fixed $d$, this statement, with the finite sequence represented
by a numerical code, is provable in iPRA.
\end{theorem}

\begin{proof}
An empty sequence causes no difficulty.  If $d=0$, a strictly descending
sequence of empty tuples contains at most one term, and $s+L\le s+1$.

For the passage from $d$ to $d+1$, partition the sequence into its
maximal consecutive blocks on which the first coordinate is constant.
The first coordinate never increases, and it strictly decreases at
every boundary between blocks.  Its initial value is at most $h(x,s)$.
Consequently there are at most $h(x,s)+1$ blocks.

Let the block boundaries, as exclusive endpoints, be
\[
 s=s_0<s_1<\cdots<s_p=s+L.
\]
Within the $i$th block the remaining $d$ coordinates form a strictly
lexicographically descending sequence, with the same coordinate envelope.
The already established $d$-dimensional theorem therefore gives
\[
                         s_{i+1}\le F_d(x,s_i).
\]
By monotonicity and induction on the number of blocks,
\[
 s_i\le F_d(x,\cdot)^{[i]}(s).
\]
Since $p\le h(x,s)+1$, this gives
\[
 s+L=s_p
 \le F_d(x,\cdot)^{[p]}(s)
 \le F_d(x,\cdot)^{[h(x,s)+1]}(s)
 =F_{d+1}(x,s).
\]

For the iPRA implementation, a block boundary is the first index at
which the first coordinate changes, found by bounded search within the
given sequence.  A block subcode is obtained by a primitive-recursive
slice operation.  The block-count inequality is proved by induction on
the bounded list of boundaries: after $i$ boundaries, the initial first
coordinate is at least the current first coordinate plus $i$.  All
conditions on sequence positions and block numbers are bounded and are
compiled into quantifier-free predicates.  No induction on an assertion
of the form ``there exists a terminating computation'' is used.
\end{proof}

The bound is exact for the stated class of chains.  This is useful for
understanding why the dimension restriction matters.  For $d=0$, use the
one-term sequence of empty tuples.  Having constructed a maximal
$d$-dimensional chain starting at any time $t$, make a $(d+1)$-dimensional
chain with $h(x,s)+1$ blocks.  In block $i$, use first coordinate
$h(x,s)-i$ and the maximal $d$-dimensional suffix chain starting at
$F_d(x,\cdot)^{[i]}(s)$.  The envelope holds because $h$ is nondecreasing;
the first coordinate drops at each block boundary; and the final
exclusive endpoint is exactly $F_{d+1}(x,s)$.  For each fixed dimension
this construction is itself primitive recursive.

\subsection{From a transition system to a bounded witness}

Let $\tau(x,u)$ be a total primitive-recursive one-step transition on
numerically coded states, and let $u_0(x)$ be the initial state.  Define
the finite simulation before proving termination:
\[
 u(x,0)=u_0(x),\qquad
 u(x,t+1)=\tau(x,u(x,t)).
\]
This is ordinary primitive recursion.  A finite simulation remains
defined even when the corresponding unbounded while-loop would fail to
halt.  Let $A(x,u)$ be a primitive-recursive activity predicate, and make
inactive states absorbing.  Let $I(x,u)$ be a quantifier-free invariant
whose initialization and preservation have iPRA proofs.  Let
\[
 r(x,u)=(r_0(x,u),\ldots,r_{d-1}(x,u))
\]
be a fixed tuple of primitive-recursive rank functions.  Assume iPRA
proves the local obligation
\begin{align*}
 I(x,u)\wedge A(x,u)\wedge A(x,\tau(x,u))
 \ \Longrightarrow\ 
 r(x,u)>_{\mathrm{lex}}r(x,\tau(x,u)).
\end{align*}

\begin{theorem}[Automatic fuel for fixed-dimensional descent]
\label{thm:automatic-lex-fuel}
Under these assumptions there is an explicitly defined
primitive-recursive $B(x)$ such that iPRA proves
\[
                         \neg A(x,u(x,B(x))).
\]
If iPRA additionally proves
$I(x,u)\wedge\neg A(x,u)\to R(x,\operatorname{out}(x,u))$,
then $w(x)=\operatorname{out}(x,u(x,B(x)))$ is a
primitive-recursive witness and iPRA proves $R(x,w(x))$.
\end{theorem}

\begin{proof}
Define the trajectory envelope
\[
 h(x,t)=\max\bigl(\{0\}\cup
 \{r_j(x,u(x,q)):q\le t,\ j<d\}\bigr).
\]
This is a bounded maximum of explicitly given primitive-recursive
values, so it is primitive recursive.  Its monotonicity and the
coordinate-envelope property are proved by quantifier-free induction on
the bounded maximum.  Set $B(x)=F_d(x,0)$ using this $h$.

Quantifier-free induction on $t$ establishes $I(x,u(x,t))$.  Suppose
that $u(x,B(x))$ were active.  Since inactive states are absorbing, all
states $u(x,0),\ldots,u(x,B(x))$ would then be active.  This finite
backward implication may be proved by contraposition of the direct
forward persistence of inactivity; activity is decidable, so this uses
only decidable reasoning.  The local rank obligation gives a controlled
strictly descending chain of $B(x)+1$ rank tuples.  Theorem
\ref{thm:controlled-lex} would imply
$B(x)+1\le F_d(x,0)=B(x)$, a contradiction.  Therefore the final state
is inactive.  The postcondition follows by applying the given local
output theorem to this state and its invariant.

There is no circularity in the envelope: $u(x,t)$ executes an already
total one-step map for an externally supplied finite number of steps.
Its definition does not inspect a termination proof.  The fixed-depth
definitions of $F_d$ then compose and iterate this particular
primitive-recursive $h$.  The resulting bound can be extremely large;
its purpose is a primitive-recursive proof of totality, not an efficient
runtime estimate.
\end{proof}

This theorem is a proof transformer, not an internal reflection principle.
For a particular transition system it consumes iPRA derivations of the
invariant and local rank conditions and produces a derivation of its
bounded termination and postcondition.  Merely proving these conditions
in Lean's ambient logic does not by itself establish their derivability
in iPRA.

\subsection{A reset example}

Consider a state $(a,b,z)\in\mathbb N^3$ and the following transition:
\[
 \tau(a,b,z)=
 \begin{cases}
 (a,b,z),&a=b=0,\\
 (a,b-1,z+1),&b>0,\\
 (a-1,2^z,z+1),&b=0<a.
 \end{cases}
\]
The active predicate is $a+b>0$.  The pair $(a,b)$ strictly decreases
lexicographically at every active step, although $a+b$ can increase
drastically at a reset.  The unused coordinate $z$ records elapsed work
and can make later resets much larger than earlier ones.

With initial data $(a_0,b_0,z_0)$, direct quantifier-free induction gives
\[
 a_t\le a_0,\qquad z_t\le z_0+t,\qquad
 b_t\le b_0+2^{z_0+t}.
\]
Thus one can use the explicit monotone envelope
$h(t)=a_0+b_0+2^{z_0+t}$ instead of the automatic maximum.  The theorem
provides primitive-recursive fuel without an induction on a termination
formula.

For this particular example one can also calculate the exact runtime.
First consume the initial $b_0$ countdown steps.  Put
\[
 T_0=b_0,\qquad T_{i+1}=T_i+1+2^{z_0+T_i}.
\]
At time $T_i$, for $i\le a_0$, the state is
$(a_0-i,0,z_0+T_i)$.  If $i<a_0$, one reset step and then exactly
$2^{z_0+T_i}$ countdown steps establish the next displayed state.
Consequently the process halts at $T_{a_0}$.  These assertions are proved
by quantifier-free inductions on the countdown length and then on $i$.
For $b_0=z_0=0$, the first runtimes as $a_0$ varies are
$0,2,7,136,2^{136}+137$.  The rapid growth is still primitive recursive:
$i\mapsto T_i$ is a fixed primitive recursion.

There cannot be a natural-valued rank that decreases under every
lexicographic decrease in $\mathbb N^2$.  Such a rank $\rho$ would satisfy
$\rho(0,b)\ge b$ and
$\rho(1,0)>\rho(0,b)$ for every $b$, which is impossible at
$b=\rho(1,0)$.  The theorem therefore does not flatten the whole order
$\omega^2$ into $\mathbb N$.  It uses the particular transition system,
its input, and its controlled finite trajectory to obtain a fuel bound.

\subsection{Why dimension must remain fixed}

The restriction is substantive, even for a simple coordinate envelope.
Take $h(t)=t+1$, suppress $x$, and write
\[
 G_0(t)=t+1,\qquad G_{d+1}(t)=G_d^{[t+2]}(t).
\]
Each fixed $G_d$ is primitive recursive, nondecreasing, and at least the
successor function.  Also $G_{d+1}(t)>G_d(t)$.

\begin{proposition}[The uniform diagonal is not primitive recursive]
The function $D(n)=G_n(n)$ is not primitive recursive.
\end{proposition}

\begin{proof}
We first prove by structural induction on primitive-recursive
definitions that, for every primitive-recursive function $f$ of any
fixed finite arity, there are fixed numbers $m,N$ such that
\[
 t\ge N,\quad x_1,\ldots,x_k\le t
 \quad\Longrightarrow\quad f(x_1,\ldots,x_k)\le G_m(t).
 \tag{$*$}
\]
Zero, successor, and projections satisfy this property.  Increase the
indices and thresholds to a common value whenever finitely many
functions are involved; this is legitimate because $G_m(t)$ is
nondecreasing in $m$ as well as in $t$.

For composition, if the component functions and the outer function are
majorized by $G_m$, the composite is majorized by
$G_m(G_m(t))\le G_{m+1}(t)$.

For primitive recursion, suppose
$f(0,\vec x)=g(\vec x)$ and
$f(n+1,\vec x)=q(n,f(n,\vec x),\vec x)$, with $g$ and $q$ majorized by
$G_m$ above a common threshold.  Fix $t$ above that threshold and
$\vec x\le t$.  Induction on $i\le t$ gives
\[
 f(i,\vec x)\le G_m^{[i+1]}(t).
\]
Indeed, the arguments of $q$ at the next step are bounded by the current
iterate, which is at least $t$.  For $n\le t$ this yields
\[
 f(n,\vec x)\le G_m^{[t+1]}(t)\le G_{m+1}(t).
\]
This proves $(*)$ for all primitive-recursive definitions.

If $D$ were primitive recursive, $(*)$ would provide fixed $m,N$ with
$D(n)\le G_m(n)$ for all sufficiently large $n$.  Choose
$n\ge\max(N,m+1)$.  Monotonicity in the level and strict growth between
successive levels give
\[
 D(n)=G_n(n)\ge G_{m+1}(n)>G_m(n),
\]
a contradiction.
\end{proof}

Consequently a proof that every fixed-dimensional instance has
primitive-recursive fuel does not furnish a primitive-recursive fuel
function whose inputs include arbitrary dimension.  In particular, an
implementation that represents the rank by a list of input-dependent
length needs an additional argument.  This distinction remains even if
Lean accepts a total definition of $(d,t)\mapsto G_d(t)$ by recursion
whose intermediate values are functions.  Lean's acceptance of such a
definition is not a certificate that it belongs to the first-order
primitive-recursive function algebra.

In fact, there is no primitive-recursive function of $(d,s)$ that bounds
the lengths of \emph{all} chains with this envelope and arbitrary
dimension.  Exactness would give
$G_d(s)\le s+B(d,s)$ for any such bound $B$.  The primitive-recursive
diagonal $n\mapsto n+B(n,n)$ would then majorize $G_n(n)$, contrary to
the same majorization argument.  Particular families with a
variable-length rank can still admit better bounds because of
additional structure; the conclusion concerns the unrestricted
lexicographic criterion.

\subsection{Lean implementation obligations}

The semantic development can use fixed tuples, a Boolean lexicographic
comparison, finite lists, and the iterator.  Its essential lemmas are:
\begin{enumerate}
 \item monotonicity and the successor lower bound for each $F_d$;
 \item the bounded block decomposition of a lexicographic chain;
 \item the controlled-chain bound, by external induction on $d$;
 \item primitive recursiveness of finite simulation, running maxima, and
       $F_d$ for each fixed $d$;
 \item preservation of a quantifier-free invariant under finite
       simulation;
 \item the contradiction from $F_d(x,0)+1$ active rank states;
 \item the output postcondition for the resulting bounded execution.
\end{enumerate}
The important typing distinction is between the defensible statement
\[
 \texttt{\textup{for every }d,\ Primrec\ (fun\ s\ =>\ fuel\ d\ s)}
\]
and the false general assertion
\[
 \texttt{Primrec\ (fun\ (d,s)\ =>\ fuel\ d\ s)}.
\]
For the object-theoretic layer, encode each fixed tuple and finite list
by natural numbers, compile each bounded assertion to a Boolean
predicate, and attach iPRA derivations of the local invariant and rank
conditions.  The fixed-dimensional construction then yields object
derivations, rather than an inference from semantic truth in Lean to
provability in iPRA.
